\documentclass[aspectratio=169,12pt]{beamer}

\usepackage[T1]{fontenc}
\usepackage[utf8]{inputenc}
\usepackage[english]{babel}
\usepackage{lmodern}
\usepackage{booktabs}
\usepackage{tabularx}

\definecolor{ink}{HTML}{19324A}
\definecolor{accent}{HTML}{176E84}
\setbeamercolor{normal text}{fg=ink,bg=white}
\setbeamercolor{frametitle}{fg=ink,bg=white}
\setbeamercolor{structure}{fg=accent}
\setbeamercolor{title}{fg=ink}
\setbeamerfont{frametitle}{size=\LARGE,series=\bfseries}
\setbeamerfont{normal text}{size=\large}
\setbeamerfont{title}{size=\Huge,series=\bfseries}
\setbeamertemplate{navigation symbols}{}
\setbeamertemplate{footline}[frame number]
\setbeamertemplate{itemize items}[circle]

\title{Discovering circuits from impedance data}
\subtitle{SOFCs, equivalent circuits, and formal grammars}
\author{Mitja Mandić}
\institute{IJS E8}
\date{}

\begin{document}

\begin{frame}[plain]
  \titlepage
\end{frame}

\begin{frame}{What is a solid oxide fuel cell?}
  \begin{itemize}
    \setlength{\itemsep}{0.95em}
    \item A SOFC converts fuel and oxygen into electricity and heat.
    \item Its hot ceramic electrolyte carries oxygen ions from the cathode to the anode.
    \item At the anode, the ions react with fuel. Released electrons flow through an external circuit.
  \end{itemize}
\end{frame}

\begin{frame}{Impedance and equivalent circuits}
  \begin{itemize}
    \setlength{\itemsep}{0.95em}
    \item EIS applies a small oscillating voltage and measures the current response across frequencies.
    \item The resulting impedance spectrum describes how the cell responds at different frequencies.
    \item An equivalent circuit combines simple elements to reproduce the spectrum.
    \item CPEs describe imperfect capacitive behavior; Warburg elements represent diffusion.
  \end{itemize}
\end{frame}

\begin{frame}{How the grammar builds circuits}
  \begin{itemize}
    \setlength{\itemsep}{0.8em}
    \item Every candidate includes a resistor in series with the rest of the circuit.
    \item Build the network by connecting allowed pieces in series or parallel.
    \item Merge equivalent resistor and inductor combinations to remove duplicates.
    \item Set a maximum element count to control the search.
  \end{itemize}
\end{frame}

\begin{frame}{The grammar choices}
  \renewcommand{\arraystretch}{1.5}
  \begin{tabularx}{\textwidth}{@{}p{0.31\textwidth}X@{}}
    \toprule
    \textbf{Grammar} & \textbf{Pieces used to build networks} \\
    \midrule
    Relaxed & Individual R, L, CPE, and W elements \\
    Compact hybrid & Predefined blocks, plus individual L and W \\
    Compact blocks & Predefined blocks only \\
    \bottomrule
  \end{tabularx}

  \vspace{0.8em}
  Example block: a resistor in parallel with a CPE.

  \vspace{0.35em}
  W denotes a Warburg element. Reported fits mainly use the relaxed collection
  without Gerischer elements.
\end{frame}

% Later sections can add parameter fitting, numerical comparisons, and results.
\end{document}
